%! Describing the Distribution of a Quantitative Variable — Unit 1, Lesson 1.4 — Lesson Plan
% GRADUAL-RELEASE plan (warm-up / guided notes and practice / debrief / close and start homework).
% THERE IS NO GROUP ACTIVITY IN THIS COURSE — the guided notes carry the release: I do / we do
% row by row in the Main Ideas / Notes table -> we do (Guided Practice, the last row). THERE IS
% NO INDEPENDENT PRACTICE SET — the homework is
% the individual practice, and the period ends with students starting it. Teacher-facing. Every box is
% `skillbox` (breakable) with a \boxguard ahead of it; this course has no `fixedskillbox`,
% no tier boxes, and no exit ticket.
\documentclass[10pt]{article}
\usepackage{apstats-article}
\usepackage{apstats-boxes}

\newcommand{\UnitNumberName}{Unit 1: Exploring One-Variable Data \quad}
\newcommand{\LessonNumberName}{Lesson 1.4: Describing the Distribution of a Quantitative Variable}

\begin{document}
\begin{center}
    {\Large\bfseries \CourseName \\
    {\small\bfseries \UnitNumberName \LessonNumberName}}
\end{center}

\begin{tcolorbox}[colback=sky, colframe=navy, boxrule=0.9pt, arc=2mm,
    left=2mm, right=2mm, top=1.5mm, bottom=1.5mm]
    \textbf{Primary Objective:} Students will describe the distribution of a quantitative
    variable by naming its shape, its unusual features --- outliers, gaps, and clusters ---
    its center, and its variability, and will state every one of those four things in the
    context of the data, naming the variable and its units.
    \\[2pt]
    \textbf{CED alignment:} Topic 1.6 \quad$\cdot$\quad Big Idea: UNC (Uncertainty)
    \quad$\cdot$\quad Skill 2.A \quad$\cdot$\quad LO UNC-1.H \quad$\cdot$\quad
    EK UNC-1.H.1--1.H.7
    \\[2pt]
    \textbf{Lesson model:} \emph{Gradual release --- I do, we do, you do --- all of it inside
    the guided notes.} There is no group activity. The notes name the vocabulary as they teach
    it and deliver the instruction row by row in one \emph{Main Ideas / Notes} table --- each
    idea a definition to complete and a grid of short problems, the first worked by you and the
    rest by the class; the last row, Guided Practice, is worked together with students holding
    the pen; \textbf{the homework is the individual practice}, and students
    start it in class while you circulate. The whole-class debrief is \emph{spoken} ---
    there is no transfer set and no reflection box on the page.
    \\[2pt]
    \textbf{Note on the ``no sketch'' rule:} every display in this lesson is \textbf{given}.
    Students construct only sentences --- they fill the pile-and-stragglers table (notes
    problem~7) and then write descriptions. Do not add a ``now draw a skewed
    distribution'' step: the AP exam asks students to \emph{read} a display and \emph{say what
    it shows}, and drawing one teaches neither.
\end{tcolorbox}

\renewcommand{\arraystretch}{1.4}
\boxguard
\begin{skillbox}[Learning Targets \& Key Understandings]{goldbox}
\begin{tabularx}{\linewidth}{@{}X|X@{}}
\textbf{Learning targets (student language)} & \textbf{Key understandings (the ``why'')} \\[2pt]
\hline
\vspace{-6pt}
\begin{itemize}[leftmargin=1.1em, nosep, after=\vspace{2pt}]
  \item I can explain why one number is never a description of a set of data.
  \item I can name a distribution's shape and say which way it leans.
  \item I can point out the values that sit apart, the empty stretches, and the separate groups.
  \item I can describe a whole distribution in four parts, in context.
\end{itemize}
&
\vspace{-6pt}
\begin{itemize}[leftmargin=1.1em, nosep, after=\vspace{2pt}]
  \item A description of a quantitative distribution states \textbf{shape, unusual features,
        center, and variability}, all four in context (EK UNC-1.H.1--1.H.6).
  \item \textbf{Target misconception:} that a distribution is named for the side the data
        \emph{piles up} on. It is named for the side the \emph{tail} runs out on. The 40 home
        prices pile on the left and the distribution is \textbf{skewed right}
        (EK UNC-1.H.3).
  \item \textbf{Second misconception:} that an outlier is an error to be deleted. It is a
        finding to be investigated; the 13-minute customer is the one the manager needs.
  \item \textbf{Third misconception:} that ``symmetric'' settles the description. The club
        heights are symmetric \emph{and} bimodal, and their center of about 68 inches belongs
        to no adult in the club.
  \item Describing these data is not a claim about a larger population --- at most it suggests
        something worth investigating (EK UNC-1.H.7).
\end{itemize}
\end{tabularx}
\end{skillbox}

\boxguard
\begin{skillbox}[{Vocabulary, Concepts, \& Theorems --- taught directly in the guided notes}]{greenbox}
\begin{minipage}[t]{0.48\linewidth}
    \begin{itemize}
        \item \textbf{Center} --- the middle, typical value (EK UNC-1.H.5)
        \item \textbf{Variability (spread)} --- how far the values reach from each other
              (EK UNC-1.H.6)
        \item \textbf{Outlier} --- a value unusually far from the rest (EK UNC-1.H.4)
        \item \textbf{Gap} --- a stretch where nothing was observed; \textbf{cluster} ---
              values bunched together, set off by a gap
    \end{itemize}
\end{minipage}
\hfill
\begin{minipage}[t]{0.48\linewidth}
    \begin{itemize}
        \item \textbf{Skewed right / skewed left} --- named for the longer tail
              (EK UNC-1.H.3)
        \item \textbf{Symmetric} --- the left half is roughly a mirror of the right
        \item \textbf{Unimodal, bimodal, uniform} --- one peak, two peaks, no peak
              (EK UNC-1.H.2)
        \item \textbf{SOCS} --- shape, outliers, center, spread; the four-part answer
        \item \textbf{Carried forward} --- dotplots and histograms, and the interval-width
              check (1.3)
    \end{itemize}
\end{minipage}
\end{skillbox}

% A tabularx never splits, so guard generously ahead of it.
\boxguard[30]
\begin{skillbox}[Lesson at a Glance --- \MeetingLength]{sky}
\renewcommand{\arraystretch}{1.25}
\begin{tabularx}{\linewidth}{@{}l c X X@{}}
\textbf{Phase} & \textbf{Min} & \textbf{Students} & \textbf{Teacher} \\
\hline
Warm-Up & 5 & Two pictures of 40 fish lengths, a middle pumpkin, and one student who read 15
  books & Debrief item~3 aloud; leave their sentence about the one student at 15 books on the
  board all period \\
Guided Notes \& Practice & 34 & Fill the vocabulary box and the notes table: problems 1--20
  row by row, then \emph{The Sleep Survey} (21--24) together, pen in their hand & \textbf{I do}
  each row's definition lines and first problem, \textbf{we do} the rest of its grid (24) $\to$
  \textbf{we do} Guided Practice (10) \\
Debrief & 8 & Pens down, papers up --- answer aloud, nothing to fill in & Put the home prices
  and the quiz scores side by side, then take the delete-the-outlier reflex and the
  symmetric-so-we're-done reflex \\
Close \& Assign & 8 & \textbf{Start the homework in class}, alone and silent & Announce the
  paper homework, circulate the first minutes for your formative read, preview Lesson~1.5 \\
\end{tabularx}
\end{skillbox}

\boxguard[20]
\begin{skillbox}[Warm-Up --- Activate Prior Knowledge (5 min)]{sky}
\begin{minipage}[t]{0.48\linewidth}
    \textbf{The three items and what each seeds}
    \begin{itemize}
        \item \textbf{1.} Forty fish lengths, two interval widths. \textbf{Spirals} to Lesson
              1.3 --- and re-establishes that a picture is a choice before today asks them to
              put confident words on one.
        \item \textbf{2.} Nine pumpkin weights in order; count in to the middle.
              \textbf{Seeds:} the center, done once with easy numbers so it is not the obstacle
              in the Center \& Variability row.
        \item \textbf{3.} A dotplot with one student far from the other fourteen.
              \textbf{Seeds:} outlier, gap, and spread, in language students already own ---
              Center \& Variability names spread and Unusual Features names the other two.
    \end{itemize}
\end{minipage}
\hfill
\begin{minipage}[t]{0.48\linewidth}
    \textbf{Running it}
    \begin{itemize}
        \item \textbf{Debrief item~3 aloud} and write a student's own phrasing on the board
              (``everyone else read two to seven books and then there's that one
              person''). Leave it up: the notes attach \emph{spread}, \emph{gap}, and
              \emph{outlier} to that sentence rather than to definitions you supply cold.
        \item Item~1's answer is \emph{neither}. If that is not automatic for most of the
              room, take ninety seconds now --- today asks students to name a shape off a
              picture, and they have to remember to check the width first.
        \item Item~2 is mechanical. Do not let it expand; nobody is computing a median today.
    \end{itemize}
\end{minipage}
\end{skillbox}

\boxguard
\begin{skillbox}[Guided Notes \& Practice --- I do, then we do (34 min)]{sky}
\textbf{This is the whole lesson.} There is no group activity, and there is no independent
practice set on this page: \textbf{the homework is the individual practice}, started in class.
The notes are one \emph{Main Ideas / Notes} table --- five rows, 24 short problems. \textbf{In
every row you fill the definition lines and work the first problem of the grid; the class works
the rest, pen in their hand.} The vocabulary box is filled \emph{as you go}, never in advance.
\begin{multicols}{2}
\textbf{Center \& Variability (5 min).} Do not rush past the setup: get a show of hands on
\emph{which line would you join} \textbf{before} anybody counts in to a middle wait, and write
the vote on the board. Work problem~1 yourself; problems 2--6 are theirs, and the middles come
out equal at problem~2 --- now the room has a problem worth five minutes. Center and spread get
named against that vote, and \textbf{problem~6 is the trap}: ``the typical wait is 4 minutes in
both lines'' is true and useless. Make them say both halves.

\textbf{Peaks \& Mirror (5 min).} The three histograms; name unimodal, bimodal, uniform, and
symmetric on the lines. Problem~7 is the pile-and-stragglers table --- fill it row by row with
the room, \emph{where does it pile, which way do the stragglers run}, because those two columns
are exactly the evidence the skew name is read off in The Tail. Do \textbf{not} name the skew
yet; problems 8--10 stay on peaks and the mirror.

\textbf{Unusual Features (5 min).} Outlier, gap, and cluster off Line B's two stranded
customers, then SOCS. Work problem~11 yourself; problem~13 is where ``in context'' stops being
a slogan --- make the rewrite name the line, the wait, and the minutes --- and problem~14 is
the one-line limit on what a description may claim.

\textbf{The Tail (9 min) --- the crux.} Count the homes live on the definition lines --- 30 of
the 40 under \$400{,}000, and the axis still runs past \$700{,}000 --- and \emph{then} say the
name, \emph{skewed right}, and let the protest happen. The sentence to land is \emph{the name
follows the tail, not the pile}; say it, write it in the callout, and point at Picture~I while
you do. Say all four for Picture~I on the lines, in context. Then the grid: \textbf{problem~15
is the crux} --- Picture~II, pile on the right, skewed \emph{left}; take a hand count before
anyone speaks. Problem~16 is the same error in a classmate's mouth. \textbf{Problems 17--19 are
the counterweights}: Picture~III is symmetric \emph{and} bimodal, ``about 68 inches'' describes
no adult in the club, and the 13-minute customer is checked, not deleted. Do not skip them.
Problem~20 is a cold call --- a skewed-right variable of their own.

\columnbreak
\textbf{We do (about 10 min): The Sleep Survey.} A dotplot in a fresh context, problems 21--24,
worked entirely together. Students hold the pen; you hold the questions: \emph{``Which way does
the tail run --- point at it.''} \quad \emph{``Is the 4-hour student a mistake?''}
\textbf{Problem~21 is the crux transferred}, and \textbf{problem~24 is the whole lesson in one
answer} --- the one to protect when time is short.

\textbf{The pen must actually change hands.} Guided Practice is the only place today where you
can watch a student decide, so do not narrate it. Put problem~21 on them cold, take a hand count
on \emph{which way does it lean} before anyone speaks, and make the room defend both answers.

Circulate and demand the reason, not the word:
\begin{itemize}[leftmargin=1.1em, nosep]
  \item \textbf{Problem~21, the crux transferred:} ``Point at the tail. Now say the name.'' A
        student who answers \emph{skewed right} is reading the pile.
  \item \textbf{Problem~22:} ``Read me the gap. Where does it start and where does it end?''
  \item \textbf{Problem~24:} ``You wrote \emph{about 8}. Eight what? Of what? Say the units.''
\end{itemize}
While circulating, pick the papers to put up in the debrief --- one problem~24 that named all
four parts in context, and one that named all four with no variable and no units attached.
\textbf{Your formative read comes twice}: here, and again in the last eight minutes as students
start the homework.
\end{multicols}
\end{skillbox}

\boxguard
\begin{skillbox}[Debrief --- whole class (8 min)]{sky}
Pens down, papers up. \textbf{This debrief is spoken} --- there is no box at the end of the
notes to fill in, so it lives entirely on the board and in the room.
\begin{multicols}{2}
\begin{enumerate}[leftmargin=1.3em, nosep, itemsep=3pt]
  \item \textbf{Open on the almost-right answer.} Somebody in problem~21 wrote
        \emph{``skewed right, because most students are on the right.''} Half of that sentence
        is a correct observation. Put Pictures I and II back up side by side and make a student
        say each name \emph{and} point at the tail that earned it. Then state it as
        EK UNC-1.H.3: the skew is named for the direction of the longer tail.
  \item Put up the two problem~24s you picked. The difference between them is not
        knowledge --- it is context. Read both aloud and let the room hear which one a stranger
        could use.
  \item Circle the warm-up sentence still on the board and take the delete-the-outlier reflex:
        \emph{``Should we throw out the student who read 15 books? What would we then not
        know?''}
  \item Take the symmetric-so-we're-done reflex last, on Picture~III: \emph{``How many adults in
        this club are 68 inches tall?''}
\end{enumerate}
\columnbreak
\textbf{Two things to demand aloud.} First, \textbf{make a student describe Picture~I in one
sentence that a reader could act on} --- skewed right, the two expensive houses, about
\$300{,}000 in the middle, prices from \$100{,}000 to \$800{,}000. Second, \textbf{make them say
what those 40 homes do \emph{not} tell you} --- nothing about houses in the next town, and
nothing about next year. Without the second half, students leave believing a description is a
conclusion.

\textbf{Connect back to Lesson~1.3 if time allows:} before you name a shape, check the interval
width that produced it. A histogram with the wrong blocks will hand you a shape that is not
there, and today you have just given students four confident words to put on it.

\textbf{If the debrief runs short}, cut the Lesson~1.3 callback and item~3 --- item~1 is the
only one that cannot be cut. \textbf{Do not let the debrief eat the last eight minutes} --- the
homework start is not optional padding, it is your second formative read.
\end{multicols}
\end{skillbox}

\boxguard
\begin{skillbox}[AP Practice --- assigned, not scored]{goldbox}
Four AP-style multiple-choice items on a 40-employee commute record --- a counting item, the
shape judgment with the ``piled on the left, so skewed left'' distractor, the over-claim item,
and the check-the-width item --- then one multi-part free response on the same display:
classify, describe with unusual features, take apart a manager's single-number summary, and
regroup at a wider interval. \textbf{Parts (a) and (d) are the spiral} --- (a) reaches back to
Lesson 1.1 (classify the variable) and Lesson 1.3 (discrete or continuous), and (d) to the
interval-width rule from Lesson 1.3. Reusing one display for both sections is deliberate: the
context is already loaded, so the thinking stays on the description.

\textbf{Not scored --- the cover's score column reads \textbf{NA} for this page.} The commute
record appears nowhere else in the packet, so what comes back in Lesson~1.5 is your evidence
that the four-part description transferred to a display students met on their own. Go over it at
the start of Lesson~1.5 and hold the AP standard out loud: \emph{in context or it does not
count}. A part (b) that says ``skewed right with an outlier'' without naming commute times and
minutes is incomplete; the answer must attach the words to the variable.
\end{skillbox}

\boxguard
\begin{skillbox}[Homework --- scored, due the first class after two study halls]{goldbox}
Five exercises spanning Topic 1.6 in fresh contexts: \textbf{1} match four displays to four
shape descriptions; \textbf{2} the crux head-on --- correct a classmate who called Display B
skewed left because the pile is on the left; \textbf{3} name the unusual feature in a bimodal
age distribution and explain it; \textbf{4} a full four-part description of 24 cereals' sugar
content; \textbf{5} an AP-style multiple-choice item with a required justification.

\textbf{Start it in the last eight minutes of the period, in class and alone.} \textbf{Scoring:}
10 points --- 2 per item. \textbf{Item~2 is where the misconception shows}: a student who
``corrects'' the classmate by agreeing that the pile decides has not learned the rule, only the
word. \textbf{Item~5 carries the check:} distractor (A) is the same error in AP clothing, and
(D) is the delete-the-outlier reflex.

\textbf{Paper homework is the default for this course} --- assign this page unless you have
decided otherwise. \textbf{DeltaMath override (occasional):} on the days you assign DeltaMath
instead, it replaces this page, you strike through the score cell on the cover, and this page
stays in the packet as extra practice. Never assign both.
\end{skillbox}

\boxguard
\begin{skillbox}[Watch For (while circulating)]{redbox}
\begin{itemize}[leftmargin=1.2em, nosep]
  \item \textbf{Skew named for the pile} (notes~15--16, GP~21, HW~1--2, AP~MC2 --- the target
        misconception). Probe: \emph{``Point at the tail. Now say the name.''}
  \item \textbf{The outlier gets deleted} (notes~19, HW~5(D), AP~MC3(D)). Probe:
        \emph{``If you throw that customer out, what do you no longer know about this line?''}
  \item \textbf{Center reported alone} (notes~6, GP~24, AP~FR(c)): the Line A / Line B trap
        surviving to the end of the period. Probe: \emph{``Both lines average four minutes.
        Which one do I join?''}
  \item \textbf{Description with no context} (notes~13, GP~24, HW~4, AP~FR(b)): ``unimodal,
        skewed right, center 6, spread 19'' with no variable and no units. Probe: \emph{``Six
        what?''}
  \item \textbf{Symmetric treated as the whole answer} (notes~17--18, HW~3): Picture~III is
        symmetric \emph{and} bimodal. Probe: \emph{``How many people are at the center?''}
  \item \textbf{Over-claiming beyond the data} (notes~14, AP~MC3, HW~5(C)): from 40 commutes to
        every worker in the city. Probe: \emph{``Who was actually measured?''}
\end{itemize}
Cold-call prompts: \emph{``Name a variable whose distribution you would expect to be skewed
right.''} \quad \emph{``Two classes have the same average. What could still be very different?''}
\end{skillbox}

\boxguard
\begin{skillbox}[Close \& Assign (8 min)]{goldbox}
\textbf{Students start the homework here, in class, alone and silent --- this is the individual
practice, and the first minutes of it are your best formative read.} Hand the launch first ---
five exercises on paper, scored, due the first class after two study halls --- and point out
that the AP Practice page before it is theirs to work and bring to review. Circulate:
\textbf{item~2 is the column that tells you whether The Tail row landed} --- a classmate has called
Display B skewed left because the pile is on the left. Sort what you see into three piles as you
walk --- said \emph{right} and pointed at the tail (ready); said \emph{right} with no reason
(ready, needs the language); kept \emph{left} or agreed with the classmate (the misconception
survived) --- and open Lesson~1.5 by reteaching The Tail row if that third pile
ran past a third of the room. Then close on what changed: \emph{``You walked in able to read a
graph. You walk out able to say what it shows --- all four parts, about a real variable, in
units, and no more than the data will carry.''} \textbf{Preview:} Lesson~1.5, \emph{Summary
Statistics for a Quantitative Variable} --- today you said ``about 8 hours'' and ``most between
6 and 9''; next time those become numbers --- mean, median, IQR, standard deviation --- and we
find out which of them one outlier can drag around and which one resists it.
\end{skillbox}

% ── Teacher notes — one per component, in packet order (four: there is no activity) ──
\begin{teachernote}[Warm-Up]
Item~3 is the seed and item~2 is the drill. Give four minutes total. \textbf{Write a student's
own phrasing of item~3 on the board verbatim and leave it up all period}; you will circle it in
the debrief when you take the delete-the-outlier reflex. Item~1's answer is \emph{neither}, and
it matters more today than it looks: this lesson hands students four confident words to put on a
shape, and the width check is the thing that stops them putting those words on an artifact.
If the room is shaky on it, spend ninety seconds and move --- do not reteach 1.3 here.
\end{teachernote}

\begin{teachernote}[Guided Notes \& Practice]
\textbf{This one document is the lesson --- pace it 24 / 10, then 8 for the debrief, then 8 for
the homework start.} The notes are a \emph{Main Ideas / Notes} table, not prose: five rows, each
a definition to complete and a grid of short problems. \textbf{In every row you fill the lines
and work the first problem; the class works the rest with the pen in their hand.} The first
three rows must move --- five minutes each. \textbf{Take the vote before the arithmetic} in
Center \& Variability: if students count in to the two middle waits first, the discovery is gone
and the row becomes a definition list; if they vote first and \emph{then} find the middles equal
at problem~2, center and variability arrive as the answer to a question they already have.
\textbf{Fill the pile-and-stragglers table (problem~7) without naming the skew.} Students will
want the word; hold it until The Tail, because the table is the evidence the name is read off,
and a name handed out first gets attached to the pile.

\textbf{Spend the time in The Tail}, which carries the misconception. Count the homes live on
the lines --- 30 of the 40 under \$400{,}000, and the axis running past \$700{,}000 --- and
\emph{then} say ``skewed right.'' Let the protest happen; it is the whole point. Do not resolve
it by repeating the definition. Point at the tail. Problem~15 is the crux --- Picture~II, pile
on the right, skewed \emph{left} --- and problem~16 is the same error in a classmate's mouth.

\textbf{Do not skip problems 17--19}, the counterweights: Picture~III, symmetric and bimodal,
center 68 inches, nobody there; and the 13-minute customer, checked rather than deleted.
Without the first, students leave believing that ``symmetric'' finishes a description, which is
the version of this error that costs points on the free response.

\textbf{The Sleep Survey is where the pen changes hands.} It is a dotplot in a fresh context,
and problem~21 is the crux transferred --- the pile is on the right and the tail runs left, so
a student reading the pile will answer \emph{skewed right}. Problem~24 is the whole lesson in
one answer. If you only get through 21, 22, and 24, that is the right trade.

\textbf{There is no independent practice set on this page, and that is deliberate --- the
homework is the individual practice.} Guided Practice is the last thing on the page; the period
ends with students starting the homework in front of you, which is where your second formative
read comes from. Item~2 is the one to watch. Sort as you circulate --- said \emph{right} and
pointed at the tail (ready); said \emph{right} with no reason (ready, needs the language); kept
\emph{left} (the misconception survived). If the third pile is more than a third of the room,
\textbf{open Lesson~1.5 by reteaching The Tail row}, and say so plainly.

\textbf{Debrief (8 min), spoken.} Open on the almost-right answer from problem~21, then
Pictures I and II side by side with a student pointing at the tail; it is the one move that
cannot be cut. \textbf{Do not let the debrief eat the last eight minutes} --- the homework start
is not optional padding.
\end{teachernote}

\begin{teachernote}[AP Practice]
\textbf{This page is not required and not scored} --- the graded practice for this lesson is the
homework behind it at the back of the packet. Work it as AP-format reps and go over it at the
start of Lesson~1.5. \textbf{MC item~2 is the one worth reading closely}: distractor (B) says
``skewed to the left, because most of the commutes lie on the left side of the graph,'' which is
the target misconception written out in full. If more than a handful pick it, spend three
minutes re-pointing at tails before starting 1.5, because this error resurfaces in Unit~1 again
at boxplots and in Unit~5 with sampling distributions. On the free response, part (c) is the
highest-value part --- it forces students to say both what a single number hides \emph{and} what
these 40 employees cannot tell you about anybody else, which is EK UNC-1.H.7 in the form the
exam actually asks it. Part (d) has a definite answer: at 30-minute intervals the bars are 27,
12, and 1, and the long commute is alone in a 60--90 bar with no empty bar in front of it.
\end{teachernote}

\begin{teachernote}[Homework]
\textbf{Scored, 10 points (2 per item), due the first class after two study halls.} The eight
minutes at the end of the period are a \emph{start}, not a completion: put students on 1, 2,
and~5 while you can still catch a wrong start, and let 3 and~4 go home. Item~2 is the one to
read first when scoring --- any student who ``corrects'' the classmate while still reasoning
from the pile has learned the word and not the rule. Item~4 is the fullest diagnostic in the
set: it is the only place a student has to produce all four parts unprompted, so score it
against the four, and dock the answer that reports numbers with no variable and no units.
\textbf{When you score the page, sort item~5 into three piles} --- chose (B) and justified it by
the tail (ready); chose (B) vaguely (ready, needs the language); chose (A), (C), or (D) (the
misconception, the over-claim, or the delete reflex survived). If more than a third land in the
last pile, reopen Pictures I and II in the Lesson~1.5 warm-up debrief rather than letting it
ride.

\textbf{Paper is the default.} This page is what gets assigned unless you decide otherwise on a
given day. On the occasions you assign DeltaMath in its place, strike the score cell on the cover
and tell students this page is now optional practice. Do not assign both.
\end{teachernote}
\end{document}
